\section{1.5 -- Im Arbeitszimmer des Präsidenten}
\noindent\textbf{Ort:} Büro des Universitätspräsidenten, Miskatonic University\\
\textbf{Zeit:} 5. April 1920, gegen 11 Uhr\\
\textbf{Anwesend:} Richard Garmond, Francis Collins, Lillie Johnson, Karim el-Mady, Natheniel Moss

\subsection{1.5.1 -- Das Arbeitszimmer}
\begin{vorlesetextalle}
Hinter der schweren Doppeltür erwartet Euch ein Raum, dessen Grö\ss{}e und Einrichtung eher an die Bibliothek eines vermögenden Gelehrten als an ein gewöhnliches Büro erinnern. Die hohe Decke ist mit feinen Stuckarbeiten versehen, und das Licht der Vormittagssonne fällt durch drei gro\ss{}e Fenster an der gegenüberliegenden Wand. Dunkle Holzvertäfelungen und der Geruch nach gewachstem Eichenholz und altem Papier verleihen dem Zimmer eine gediegene, beinahe ehrfurchtgebietende Atmosphäre.

Ein kostbarer, dicht geknüpfter Teppich in tiefen Rot-, Blau- und Goldtönen bedeckt einen beträchtlichen Teil des Parkettbodens. Seine aufwendigen, ineinanderlaufenden Muster wirken, als habe man sie aus einem fernen Land hierhergebracht. Darüber hinaus dämpft er jeden Schritt so vollständig, dass sich beim Betreten des Raumes unwillkürlich das Gefühl einstellt, die Geräusche des übrigen Universitätsbetriebs hinter sich gelassen zu haben.

Nahe den Fenstern steht ein mächtiger Schreibtisch aus dunklem Holz, dessen breite Platte von einer grünen Schreibunterlage, sorgfältig aufgeschichteten Schriftstücken, mehreren Büchern und einer schweren Messinglampe eingenommen wird. Dahinter sitzt ein Herr in einem tadellos geschnittenen dunklen Anzug, der gerade über einem Dokument gebeugt ist.

Die Wände sind fast vollständig von raumhohen, fest eingebauten Bücherregalen bedeckt. Hinter zum Teil verglasten Türen reihen sich Hunderte, vielleicht Tausende Bände aneinander: alte Folianten in Leder, wissenschaftliche Reihen in schlichten Leineneinbänden und jüngere Veröffentlichungen mit fein geprägten Rücken. Kleine Messingleitern stehen bereit, um auch die oberen Reihen zu erreichen. Hier und da unterbrechen antike Büsten und sorgfältig beschriftete Fundstücke die langen Reihen von Büchern.

Seitlich vor dem Schreibtisch ist eine Sitzecke eingerichtet. Vier tiefe, weich gepolsterte Ledersessel stehen um ein niedriges Tischchen, auf dem drei Gläser und eine Karaffe mit bernsteinfarbener Flüssigkeit bereitstehen. Ein wenig abseits des Ensembles befindet sich ein einfacher Holzstuhl: gerade Lehne, keine Polsterung, sichtbar weniger bequem als die übrigen Plätze. Er wirkt nicht beschädigt oder zufällig herbeigeschafft, sondern schlicht so, als gehöre es sich hier, dass nicht jeder Gast denselben Sitzplatz erhält.

Der Herr hinter dem Schreibtisch hebt den Blick von seinen Unterlagen.
\end{vorlesetextalle}

\begin{meisterinfo}
\textbf{Probe: Bildung oder Geschichte (alternativ Bibliotheksnutzung).} Sobald die Ermittler Gelegenheit haben, sich die Bücherregale anzusehen, darf jede interessierte Person eine passende Probe ablegen. Der Hinweis darf auch ohne Probe durch gezieltes Nachfragen erlangt werden; der Wurf vermittelt ihn nur frühzeitig durch Beobachtung.

\textbf{Erfolg:} Zwischen den Beständen fallen grö\ss{}ere Reihen zur Alten Geschichte, vergleichenden Kulturgeschichte und Archäologie der Mittelmeerwelt auf. Einige Veröffentlichungen tragen den Namen \emph{Richard Garmond}. Offenbar ist der Universitätspräsident nicht nur Verwalter, sondern selbst ein ausgewiesener Gelehrter auf dem Gebiet der Alten Geschichte und Archäologie; ein besonderes Fachinteresse an Ägypten lässt sich daraus nicht ableiten.

\textbf{Misserfolg:} Die Bibliothek ist erkennbar umfangreich und wissenschaftlich anspruchsvoll, aber aus der kurzen Betrachtung lässt sich kein eindeutiger fachlicher Schwerpunkt bestimmen.

\textbf{Gesellschaftlicher Unterschied:} Der einfache Stuhl ist für Natheniel Moss, der als Beschäftigter der Universität an der Besprechung teilnimmt. Dies ist keine gezielte Demütigung: Garmond und seine Umgebung nehmen die Rangordnung als selbstverständlich wahr. Die Spielenden dürfen diese Geste selbstverständlich wahrnehmen und darauf reagieren. Die drei Gläser stehen zunächst für die externen Gäste bereit; Natheniel wird nichts angeboten, es sei denn, jemand spricht es an.
\end{meisterinfo}

\subsection{1.5.2 -- Begrü\ss{}ung und Sitzordnung}
\begin{vorlesetextalle}
Der Mann schiebt das Schriftstück, das er eben gelesen hat, beiseite, erhebt sich hinter seinem Schreibtisch und kommt Euch mit gemessenen Schritten entgegen. Er ist etwa fünfzig Jahre alt, seine Schläfen sind bereits leicht ergraut. Seine Haltung ist aufrecht und selbstsicher, doch sein Lächeln wirkt offen und aufrichtig.

\glqq{}Meine Damen, meine Herren, herzlich willkommen. Richard Garmond. Es freut mich, dass Sie alle meiner Einladung gefolgt sind.\grqq{}

Dann fällt sein Blick auf Francis Collins, und sein Gesicht hellt sich sichtbar auf.

\glqq{}Francis! Wie schön, Dich wiederzusehen. Nach Deinen Unternehmungen im vergangenen Jahr habe ich gehofft, Dich bald wieder hier begrü\ss{}en zu dürfen. Ich hoffe, die Reise war nicht allzu beschwerlich.\grqq{}

Er wendet sich den übrigen Gästen zu und begrü\ss{}t sie der Reihe nach mit einem Händedruck.

\glqq{}Fräulein Johnson, vielen Dank, dass Sie gekommen sind. Herr el-Mady, ich freue mich, Sie persönlich kennenzulernen.\grqq{}

Auch Natheniel Moss erhält einen kurzen, höflichen Gru\ss{}, doch Garmonds Aufmerksamkeit richtet sich sogleich wieder auf die anderen.

\glqq{}Aber bitte, nehmen Sie doch Platz. Wir haben einiges zu besprechen.\grqq{}

Mit einer einladenden Geste weist er auf die gepolsterten Sessel. Sein Blick streift den einfachen Holzstuhl, als sei dessen Bestimmung ohnehin klar.
\end{vorlesetextalle}

\begin{meisterinfo}
\textbf{Beziehungen:} Garmond kennt Francis Collins persönlich und duzt sie; er spricht Lillie Johnson und Karim el-Mady mit \emph{Sie} an. Natheniel Moss ist Universitätsangestellter und wird korrekt, aber eher beiläufig behandelt.

\textbf{Sitzordnung:} Die vier Polstersessel sind als repräsentative Gästeplätze arrangiert. Der Holzstuhl ist für Moss vorgesehen; der Präsident erklärt die Hierarchie nicht ausdrücklich. Sollte jemand einen anderen Platz anbieten oder die Ungleichbehandlung ansprechen, reagiert Garmond eher überrascht als verärgert und lässt sich auf eine unkomplizierte Lösung ein.

\textbf{Getränke:} Auf dem Tisch stehen nur drei Gläser. Der Präsident kann den Gästen einen Whiskey anbieten; ein zusätzliches Glas kann bei Nachfrage gebracht werden. Niemand muss etwas trinken.
\end{meisterinfo}

\subsection{1.5.3 -- Der Auftrag}
\begin{vorlesetextalle}
Nachdem Ihr Platz genommen habt, bleibt Professor Garmond einen Moment stehen, als suche er nach den richtigen Worten. Schlie\ss{}lich setzt er sich zu Euch und legt die Fingerspitzen gegeneinander.

\glqq{}Ich möchte Ihnen zunächst danken, dass Sie sich heute die Zeit genommen haben. Der Anlass unserer Zusammenkunft ist eine Angelegenheit, die mich persönlich und die Universität zunehmend beunruhigt.\grqq{}

Er blickt kurz in die Runde.

\glqq{}Im vergangenen Herbst brach eine von der Miskatonic University unterstützte Expedition nach Ägypten auf. Sie stand unter der Leitung von Dr. Skorbinsk, einem erfahrenen Wissenschaftler, der uns auf eine bisher kaum erforschte Fundstätte südwestlich von Luxor aufmerksam gemacht hatte. Begleitet wurde er unter anderem von Flynn Carlson, den einige von Ihnen möglicherweise kennen.\grqq{}

Garmond nimmt eines der Papiere vom Tisch, betrachtet es flüchtig und legt es wieder ab.

\glqq{}Die ersten Berichte waren au\ss{}erordentlich vielversprechend. Unsere Leute hatten Überreste einer alten Anlage entdeckt, deren Umfang und Beschaffenheit sich zunächst kaum abschätzen lie\ss{}en. Es bestand durchaus Anlass zu der Hoffnung, dass wir es mit einer bemerkenswerten archäologischen Entdeckung zu tun hatten.\grqq{}

Sein Gesicht wird ernster.

\glqq{}Doch dann wurden die Nachrichten immer spärlicher. Wir hörten von Schwierigkeiten mit den angeworbenen Arbeitern, von Verzögerungen und von einer Unruhe, für die uns niemand eine überzeugende Erklärung geben konnte. Schlie\ss{}lich brach der unregelmä\ss{}ig gewordene Kontakt ganz ab.\grqq{}

Er lässt eine kurze Pause entstehen.

\glqq{}Eine letzte Nachricht erreichte uns über einen Mann namens Aziz, der für die Expedition als Übersetzer tätig war. Er meldete sich aus Luxor und berichtete, dass die Situation am Grabungsort bedenklich sei. Seine Angaben waren bruchstückhaft. Seitdem haben wir keine verlässliche Auskunft mehr über den Verbleib von Dr. Skorbinsk und Flynn Carlson.\grqq{}

Garmond lehnt sich vor.

\glqq{}Sie werden verstehen, dass ich die Angelegenheit nicht länger auf sich beruhen lassen kann. Wir tragen die Verantwortung für unsere Mitarbeiter und für die Menschen, die wir in unserem Namen in die Wüste geschickt haben. Zugleich müssen wir feststellen, was aus den bisherigen Forschungsergebnissen geworden ist.\grqq{}

Er sieht nacheinander Francis, Lillie und Karim an.

\glqq{}Daher möchte ich Sie bitten, als kleine, unabhängige Gruppe nach Ägypten zu reisen. Suchen Sie zunächst in Kairo und Luxor nach weiteren Nachrichten, stellen Sie nach Möglichkeit Verbindung zu Aziz her und begeben Sie sich anschlie\ss{}end zur Ausgrabungsstätte. Finden Sie heraus, was dort geschehen ist und ob unsere Leute Hilfe benötigen. Sollte die Lage es erlauben, möchte ich au\ss{}erdem eine nüchterne Einschätzung der bisherigen Funde erhalten.\grqq{}

Er deutet kurz zu Natheniel hinüber.

\glqq{}Herr Moss wird Sie dabei mit seinen Kenntnissen der alten Schriften und der ägyptischen Frühgeschichte unterstützen. Die Universität hat seine Teilnahme bereits veranlasst.\grqq{}

Dann wendet er sich wieder an die übrigen.

\glqq{}Die Universität übernimmt die notwendigen Kosten der Reise, Ihrer Unterbringung und Ausrüstung. Die Einzelheiten der Vergütung wurden Ihnen bereits schriftlich mitgeteilt. Ich hoffe sehr, dass ich auf Ihre Unterstützung zählen darf.\grqq{}
\end{vorlesetextalle}

\begin{meisterinfo}
\textbf{Pause für Fragen und Zusage:} Halte den Monolog an dieser Stelle an. Die Spieler dürfen nach dem genauen Auftrag, der letzten Nachricht von Aziz, der Reise, der Finanzierung oder der Dauer fragen. Garmond wei\ss{}, dass die Expedition südwestlich von Luxor tätig war; er hat keine gesicherte Erklärung für das Schweigen der Forscher. Er stellt keine übernatürlichen Vermutungen an.

\textbf{Wichtig:} Die individuelle Vergütung ergibt sich aus den bereits ausgeteilten Einladungsbriefen. Für Natheniel ist die Reise eine dienstliche Anweisung, kein neues freiwilliges Arbeitsangebot. Garmond erwartet seine Mitwirkung, ohne ihn ständig persönlich anzusprechen.

\textbf{Falls die Gruppe nicht sofort zusagt:} Garmond bittet um eine baldige Entscheidung und beantwortet praktische Fragen. Die Handlung sollte nicht durch eine erzwungene Zustimmung blockiert werden.
\end{meisterinfo}

\subsection{1.5.4 -- Der Hinweis auf das Carlson-Haus}
\begin{vorlesetextalle}
Als die wesentlichen Fragen besprochen sind, hält Professor Garmond noch einmal inne, als sei ihm ein letzter Gedanke gekommen.

\glqq{}Bevor Sie jedoch Ihre Reise antreten, möchte ich Ihnen noch etwas ans Herz legen. Flynn Carlson hat seiner Frau während der Expedition offenbar regelmä\ss{}ig geschrieben. Mrs. Carlson lebt hier in Arkham, nur wenige Meilen von der Universität entfernt.\grqq{}

Er nimmt seine Brille ab und legt sie neben die Unterlagen.

\glqq{}Möglicherweise hat sie Briefe erhalten, die uns weitere Aufschlüsse über die letzten Wochen der Expedition geben können. Ich würde Ihnen daher dringend empfehlen, ihr einen Besuch abzustatten. Vielleicht erfahren Sie dort etwas, das Ihnen die Suche in Ägypten erleichtert.\grqq{}

Der Präsident sieht Euch erwartungsvoll an.

\glqq{}Haben Sie noch Fragen, bevor wir die weiteren Vorbereitungen treffen?\grqq{}
\end{vorlesetextalle}

\begin{meisterinfo}
\textbf{Übergang zu Szene 2:} Garmond gibt der Gruppe auf Nachfrage den bekannten Wohnort der Carlsons in Arkham und ermöglicht ihr den Besuch. Das ist der ausdrücklich in den älteren Meister-Sheets vorgesehene Handlungsimpuls. Die Spieler dürfen vor der Reise auch andere Vorbereitungen treffen; der Besuch des Hauses wird jedoch klar empfohlen.

\textbf{Szenenende:} Wenn die Ermittler das Büro verlassen, ist Szene 1 abgeschlossen. Es folgt \emph{Szene 2 -- Das verlassene Haus der Carlsons}.
\end{meisterinfo}
